\documentclass[12pt, titlepage]{article}

\usepackage{amsmath, mathtools}

\usepackage[round]{natbib}
\usepackage{amsfonts}
\usepackage{amssymb}
\usepackage{graphicx}
\usepackage{colortbl}
\usepackage{xr}
\usepackage{hyperref}
\usepackage{longtable}
\usepackage{xfrac}
\usepackage{tabularx}
\usepackage{float}
\usepackage{siunitx}
\usepackage{booktabs}
\usepackage{multirow}
\usepackage[section]{placeins}
\usepackage{caption}
\usepackage{fullpage}

\hypersetup{
bookmarks=true,     % show bookmarks bar?
colorlinks=true,       % false: boxed links; true: colored links
linkcolor=red,          % color of internal links (change box color with linkbordercolor)
citecolor=blue,      % color of links to bibliography
filecolor=magenta,  % color of file links
urlcolor=cyan          % color of external links
}

\usepackage{array}

\externaldocument{../../SRS/SRS}

\input{../../Comments.text}
\input{../../Common.text}

\begin{document}

\title{Module Interface Specification for \progname{}}

\author{\authname}

\date{\today}

\maketitle

\pagenumbering{roman}

\section{Revision History}

\begin{tabularx}{\textwidth}{p{3cm}p{2cm}X}
\toprule {\bf Date} & {\bf Version} & {\bf Notes}\\
\midrule
March 20, 2026 & 1.0 & Initial version\\
April 21 & 1.1 & Update according to smith's feedback\\
\bottomrule
\end{tabularx}

~\newpage

\section{Symbols, Abbreviations and Acronyms}

See SRS Documentation at \href{https://github.com/Krozix-2026/cas741-speech-trf/blob/main/docs/SRS/SRS.tex}{SRS.tex}



% \wss{Also add any additional symbols, abbreviations or acronyms}

\newpage

\tableofcontents

\newpage

\pagenumbering{arabic}

\section{Introduction}

The following document details the Module Interface Specifications for
\progname{}. This pipeline accepts different deep learning-based models for MEG analyzing, performing mTRF test for comparison.


Complementary documents include the System Requirement Specifications
and Module Guide.  The full documentation and implementation can be
found at \url{https://github.com/Krozix-2026/cas741-speech-trf}.  
The key companion documents are \href{https://github.com/Krozix-2026/cas741-speech-trf/blob/main/docs/SRS/SRS.tex}{SRS}
and
\href{https://github.com/Krozix-2026/cas741-speech-trf/blob/main/docs/Design/SoftArchitecture/MG.tex}{MG}.

\section{Notation}

% \wss{You should describe your notation.  You can use what is below as
%   a starting point.}

The structure of the MIS for modules comes from \citet{HoffmanAndStrooper1995},
with the addition that template modules have been adapted from
\cite{GhezziEtAl2003}.  The mathematical notation comes from Chapter 3 of
\citet{HoffmanAndStrooper1995}.  
% For instance, the symbol := is used for a
% multiple assignment statement and conditional rules follow the form $(c_1
% \Rightarrow r_1 | c_2 \Rightarrow r_2 | ... | c_n \Rightarrow r_n )$.

The following table summarizes the primitive data types used by \progname. 

\renewcommand{\arraystretch}{1.15}
\noindent
\begin{longtable}{l l p{11.0cm}}
\toprule
\textbf{Symbol} & \textbf{Unit} & \textbf{Description}\\
\midrule
\endfirsthead

\toprule
\textbf{Symbol} & \textbf{Unit} & \textbf{Description}\\
\midrule
\endhead

\midrule
\multicolumn{3}{r}{\emph{Continued on next page}}\\
\endfoot

\bottomrule
\endlastfoot

$C$ & -- & Number of MEG sensors.\\
$D$ & -- & Predictor feature dimension.\\
$D_b$ & -- & Feature dimension of the baseline predictor matrix $\mathbf{R}^{(\mathrm{base})}$.\\
$B$ & -- & Training batch size.\\
$E$ & -- & Number of training epochs.\\
$K$ & -- & Number of time lags used in the TRF model.\\
$L$ & -- & Number of candidate predictors or layers being compared.\\
$N$ & -- & Number of aligned time samples used for TRF fitting.\\
$N_{\mathrm{ROI}}$ & -- & Number of regions of interest (ROI), if ROI-based analysis is performed.\\

$\mathcal{D}$ & -- & Speech dataset collection.\\
$\mathcal{D}_{\mathrm{train}}$ & -- & Training split of dataset $\mathcal{D}$.\\
$\mathcal{D}_{\mathrm{val}}$ & -- & Validation split of dataset $\mathcal{D}$.\\
$\mathcal{D}_{\mathrm{test}}$ & -- & Test split of dataset $\mathcal{D}$.\\

$\boldsymbol{a}$ & -- & Audio input sequence provided to the speech model.\\
$\boldsymbol{s}^{\star}$ & -- & Ground-truth transcript for ASR training.\\
$\boldsymbol{\theta}$ & -- & Trainable parameters of the speech model.\\
$\boldsymbol{\theta}^{\ast}$ & -- & Optimized model parameters after training.\\
$f_{\boldsymbol{\theta}}(\cdot)$ & -- & Speech model parameterized by $\boldsymbol{\theta}$.\\
$\mathcal{L}_{\mathrm{ASR}}$ & -- & ASR training loss used to fit $f_{\boldsymbol{\theta}}$.\\
$\eta$ & -- & Learning rate for gradient-based optimization.\\

$\ell$ & -- & Layer index of the deep speech model.\\
$H_{\ell}$ & -- & Hidden-state dimensionality of layer $\ell$.\\
$n$ & -- & Discrete time index after temporal alignment or model frame sampling.\\
$t$ & \si{\second} & Continuous time variable in seconds.\\
$\Delta t$ & \si{\second} & Sampling interval of MEG signals, $\Delta t = 1/f_s$.\\

$\boldsymbol{h}^{(\ell)}_{n}$ & -- & Hidden state vector from layer $\ell$ at discrete time index $n$.\\
$g(\cdot)$ & -- & Optional transformation applied to hidden states to form predictors.\\
$\boldsymbol{r}^{(\ell)}_{n}$ & -- & Predictor representation vector derived from $\boldsymbol{h}^{(\ell)}_{n}$.\\
$\mathbf{R}^{(\ell)}$ & -- & Time-by-feature predictor matrix derived from layer $\ell$, with $\mathbf{R}^{(\ell)} \in \mathbb{R}^{N \times D}$.\\
$\mathbf{R}^{(\mathrm{base})}$ & -- & Baseline predictor matrix, with $\mathbf{R}^{(\mathrm{base})} \in \mathbb{R}^{N \times D_b}$.\\

$\mathbf{X}(\mathbf{R})$ & -- & Lagged design matrix constructed from predictor matrix $\mathbf{R}$, with $\mathbf{X}(\mathbf{R}) \in \mathbb{R}^{N \times (KD)}$.\\
$\mathbf{X}$ & -- & Design matrix formed from time-lagged predictors when the underlying predictor matrix is clear from context.\\
$\mathbf{X}_{\tau}$ & -- & Portion of the design matrix corresponding to lag $\tau$.\\

$f_s$ & \si{\hertz} & Sampling frequency of MEG signals used for TRF modeling.\\
$f_a$ & \si{\hertz} & Sampling frequency of the raw audio waveform.\\
$\tau$ & \si{\second} & Time-lag variable in the TRF model.\\
$\tau_{\min}$ & \si{\second} & Start of the TRF lag window.\\
$\tau_{\max}$ & \si{\second} & End of the TRF lag window.\\

$\boldsymbol{x}(n)$ & -- & Predictor vector at aligned discrete time index $n$.\\
$\boldsymbol{y}(n)$ & \si{\tesla} / \si{\femto\tesla} & MEG response vector across sensors at aligned discrete time index $n$.\\
$\mathbf{Y}$ & \si{\tesla} / \si{\femto\tesla} & MEG response matrix across time and sensors.\\
$\hat{\mathbf{Y}}$ & \si{\tesla} / \si{\femto\tesla} & Predicted MEG responses from the fitted encoding model.\\

$\mathbf{w}(\tau)$ & -- & TRF weight vector at lag $\tau$.\\
$\mathbf{W}$ & -- & Full TRF weight matrix across lags and sensors.\\
$\mathbf{W}^{\ast}$ & -- & Ridge-regularized estimate of the TRF weights.\\
$\lambda$ & -- & Ridge regularization coefficient for TRF fitting.\\

$\rho$ & -- & Prediction score of an encoding model.\\
$\rho_{\mathrm{base}}$ & -- & Encoding score obtained using only the baseline predictor.\\
$\rho_{\ell}$ & -- & Encoding score obtained using predictors derived from layer $\ell$.\\
$\Delta \rho_{\ell}$ & -- & Improvement over baseline, defined as $\Delta \rho_{\ell} = \rho_{\ell} - \rho_{\mathrm{base}}$.\\
$\rho_{\mathrm{CV}}$ & -- & Cross-validated prediction score.\\

\end{longtable}


\noindent
The specification of \progname \ uses some derived data types: sequences,
strings, tuples, matrices, and tensors. Sequences are lists filled with
elements of the same data type. Strings are sequences of characters. Tuples
contain a list of values, potentially of different types. Matrices and tensors
denote multi-dimensional arrays of real values. In addition, \progname \ pipeline uses
functions, which are defined by the data types of their inputs and outputs.
Local functions are described by giving their type signature followed by their
specification.

\section{Module Decomposition}

The following table is taken directly from the Module Guide document for this project.

\begin{table}[h!]
\centering
\begin{tabular}{p{0.3\textwidth} p{0.6\textwidth}}
\toprule
\textbf{Level 1} & \textbf{Level 2}\\
\midrule

\multirow{2}{0.3\textwidth}{Hardware-Hiding}
& Audio Data Module\\
& MEG Data Module\\
\midrule

\multirow{5}{0.3\textwidth}{Behaviour-Hiding}
& Input Parameters Module\\
& Model Training Module\\
& Frozen Model Module\\
& Representation Extraction Module\\
% & Baseline Predictor Module\\
% & Predictor Assembly Module\\
% & Brain Alignment Module\\
& Statistical Analysis Module\\
% & Output Verification Module\\
% & Output Format Module\\
\midrule

\multirow{1}{0.3\textwidth}{Software Decision}
% & Checkpoint / Serialization Module\\
% & Feature Cache Module\\
% & TRF Backend Interface Module\\
& Plotting Module\\
\bottomrule
\end{tabular}
\caption{Module Hierarchy}
\label{TblMH}
\end{table}







\newpage
% ~\newpage



% \section{MIS of \wss{Module Name}} \label{Module} \wss{Use labels for
%   cross-referencing}

% \wss{You can reference SRS labels, such as R\ref{R_Inputs}.}

% \wss{It is also possible to use \LaTeX for hypperlinks to external documents.}

% \subsection{Module}

% \wss{Short name for the module}

% \subsection{Uses}


% \subsection{Syntax}

% \subsubsection{Exported Constants}

% \subsubsection{Exported Access Programs}

% \begin{center}
% \begin{tabular}{p{2cm} p{4cm} p{4cm} p{2cm}}
% \hline
% \textbf{Name} & \textbf{In} & \textbf{Out} & \textbf{Exceptions} \\
% \hline
% \wss{accessProg} & - & - & - \\
% \hline
% \end{tabular}
% \end{center}

% \subsection{Semantics}

% \subsubsection{State Variables}

% \wss{Not all modules will have state variables.  State variables give the module
%   a memory.}

% \subsubsection{Environment Variables}

% \wss{This section is not necessary for all modules.  Its purpose is to capture
%   when the module has external interaction with the environment, such as for a
%   device driver, screen interface, keyboard, file, etc.}

% \subsubsection{Assumptions}

% \wss{Try to minimize assumptions and anticipate programmer errors via
%   exceptions, but for practical purposes assumptions are sometimes appropriate.}

% \subsubsection{Access Routine Semantics}

% \noindent \wss{accessProg}():
% \begin{itemize}
% \item transition: \wss{if appropriate} 
% \item output: \wss{if appropriate} 
% \item exception: \wss{if appropriate} 
% \end{itemize}

% \wss{A module without environment variables or state variables is unlikely to
%   have a state transition.  In this case a state transition can only occur if
%   the module is changing the state of another module.}

% \wss{Modules rarely have both a transition and an output.  In most cases you
%   will have one or the other.}

% \subsubsection{Local Functions}

% \wss{As appropriate} \wss{These functions are for the purpose of specification.
%   They are not necessarily something that is going to be implemented
%   explicitly.  Even if they are implemented, they are not exported; they only
%   have local scope.}


\section{MIS of Audio Data Module} \label{Module_AudioData}

\subsection{Module}

Audio Data Module

\subsection{Uses}

Input Parameters Module

\subsection{Syntax}

\subsubsection{Exported Constants}

\begin{itemize}
\item $\mathcal{S}_{\mathrm{Audio}}$: the finite set of supported audio dataset identifiers. In the current implementation this includes at least the LibriSpeech-based training data used by the speech-model training pipeline.

\item $\mathcal{F}_{\mathrm{Seq}}$: the abstract type of a time-by-feature
matrix, represented in the current implementation as a tensor of shape
$(T, F)$.

\item $\mathcal{B}_{\mathrm{Audio}}$: the abstract type of a batched audio
training sample. In the current implementation this is represented by
\texttt{LibriSpeechBatch}, \texttt{AlignedWordBatch}, or an equivalent
dictionary-based batch record, depending on the training task.
\end{itemize}

\subsubsection{Exported Access Programs}

\begin{center}
\begin{tabular}{p{4.6cm} p{4.6cm} p{3.8cm} p{3.0cm}}
\hline
\textbf{Name} & \textbf{In} & \textbf{Out} & \textbf{Exceptions} \\
\hline
LibriSpeechASR &
root: Path, subset: string, vocab: CharVocab, sample\_rate: int, n\_mels: int, win\_length\_ms: real, hop\_length\_ms: real, download: bool, limit\_samples: optional int &
dataset object &
dataset unavailable, metadata unavailable, file error \\

LibriSpeechASR.\\\_\_getitem\_\_ &
idx: int &
record containing feat, feat\_len, target, target\_len, text, utt\_id &
index out of range, file error, malformed audio \\

LibriSpeechAlignedWords &
manifest\_path: Path, subset: string, vocab: WordVocab, tail\_frames: int, limit\_samples: optional int &
dataset object &
manifest unavailable, malformed manifest, feature file error \\

LibriSpeechAlignedWords.\\\_\_getitem\_\_ &
idx: int &
record containing feat, feat\_len, target, word\_ids, word\_starts, word\_ends, utt\_id &
index out of range, malformed feature array, file error \\

buildWordVocabFrom\\Manifest &
manifest\_path: Path, topk: int &
WordVocab &
manifest unavailable, malformed manifest \\

collateLibriSpeech &
sequence of LibriSpeechASR item records &
LibriSpeechBatch &
malformed item record, inconsistent feature dimension \\

collateRNNT &
sequence of LibriSpeechASR item records &
dictionary batch record &
malformed item record, inconsistent feature dimension \\

collateAlignedWords &
sequence of LibriSpeechAlignedWords item records &
AlignedWordBatch &
malformed item record, inconsistent feature dimension \\
\hline
\end{tabular}
\end{center}

\subsection{Semantics}

\subsubsection{State Variables}

\begin{itemize}
\item $subset$: string, the dataset subset currently represented by a dataset object (for example, train, validation, or test subset).

\item $indices$: optional sequence of integers, the restricted sample indices used when \texttt{limit\_samples} is specified.

\item $sampleRate$: positive integer, the target sampling rate used for audio loading and optional resampling.

\item $featureSpec$: record containing the feature-extraction parameters used by the audio dataset object, including $n\_mels$, $win\_length\_ms$, and $hop\_length\_ms$.

\item $manifestItems$: sequence of manifest records used by
\texttt{LibriSpeechAlignedWords}. Each manifest record contains at least the subset identifier, the cochlear-feature path, the word alignment list, and the utterance identifier.

\item $vocab$: the vocabulary object used to encode targets. Depending on the task, this is either a character vocabulary or a word vocabulary.
\end{itemize}

\subsubsection{Environment Variables}

\begin{itemize}
    \item Audio files, feature files, and manifest files stored in the external file system.
\end{itemize}

\subsubsection{Assumptions}

\begin{itemize}
    \item The dataset paths and manifest paths supplied by the Input Parameters Module are valid.

    \item For \texttt{LibriSpeechASR}, each referenced audio file can be decoded into a waveform and, if necessary, resampled to the requested target sampling rate.

    \item For \texttt{LibriSpeechAlignedWords}, each manifest record contains a valid \texttt{coch\_path} whose array has shape $(64, T)$.

    \item Each utterance identifier is unique within the dataset subset being used.
\end{itemize}

\subsubsection{Access Routine Semantics}

\noindent \textbf{LibriSpeechASR}(root, subset, vocab, ...):
\begin{itemize}
    \item transition: initializes a dataset object backed by
    \texttt{torchaudio.datasets.LIBRISPEECH}, stores the requested subset and feature-extraction parameters, and prepares the mel-spectrogram and amplitude-to-dB transforms.
    \item exception: raised if the requested subset cannot be loaded or if the dataset root is invalid.
\end{itemize}

\noindent \textbf{LibriSpeechASR.\_\_getitem\_\_}(idx):
\begin{itemize}
    \item output: uses dataset metadata to locate the corresponding audio file, loads the waveform through the local audio-loading routine, optionally resamples it to the configured sampling rate, computes the log-mel feature matrix of shape $(T, F)$, normalizes the transcript text, encodes the target sequence, and returns a record containing the feature matrix, feature length, target sequence, target length, normalized text, and utterance identifier.
    \item exception: raised if the index is invalid, the audio file cannot be read, or the waveform cannot be transformed into the required feature format.
\end{itemize}

\noindent \textbf{LibriSpeechAlignedWords}(manifest\_path, subset, vocab, ...):
\begin{itemize}
    \item transition: loads the manifest records for the requested subset,
    stores the tail-frame setting and the vocabulary object, and prepares the aligned-word dataset object.
    \item exception: raised if the manifest file is unavailable or malformed.
\end{itemize}

\noindent \textbf{LibriSpeechAlignedWords.\_\_getitem\_\_}(idx):
\begin{itemize}
    \item output: loads the cochlear feature array referenced by the selected manifest record, converts it into a time-by-feature tensor, constructs the frame-level target sequence together with the word-id, word-start, and word-end sequences, and returns the corresponding item record.
    \item exception: raised if the index is invalid, the feature file is missing,
    or the loaded feature array has an unexpected shape.
\end{itemize}

\noindent \textbf{buildWordVocabFromManifest}(manifest\_path, topk):
\begin{itemize}
    \item output: scans the manifest file, counts word occurrences, and returns a vocabulary object containing the most frequent \texttt{topk} word types together with an unknown-token identifier.
    \item exception: raised if the manifest file cannot be read or contains malformed records.
\end{itemize}

\noindent \textbf{collateLibriSpeech}(items):
\begin{itemize}
    \item output: sorts item records by feature length, pads the feature tensors to a common time dimension, concatenates the target sequences, and returns a \texttt{LibriSpeechBatch}.
    \item exception: raised if the item records are malformed or have inconsistent feature dimensions.
\end{itemize}

\noindent \textbf{collateRNNT}(items):
\begin{itemize}
    \item output: sorts item records by feature length, pads feature tensors and target tensors separately, and returns a dictionary-based batch record for RNNT-style training.
    \item exception: raised if the item records are malformed or have inconsistent feature dimensions.
\end{itemize}

\noindent \textbf{collateAlignedWords}(items):
\begin{itemize}
    \item output: sorts aligned-word item records by feature length, pads the feature tensors, target tensors, and word-boundary arrays, and returns an \texttt{AlignedWordBatch}.
    \item exception: raised if the item records are malformed or have inconsistent feature dimensions.
\end{itemize}

\subsubsection{Local Functions}

\begin{itemize}
    \item \textbf{\_load\_audio\_soundfile}: loads an audio waveform from disk and converts stereo input to mono when needed.

    \item \textbf{normalize\_text}: converts transcripts into the normalized text form used for target encoding.

    \item \textbf{feature extraction transform}: converts waveforms into log-mel features using the configured mel-spectrogram and amplitude-to-dB transforms.
\end{itemize}












\newpage


\section{MIS of MEG Data Module} \label{Module_MEGData}

\subsection{Module}

MEG Data Module

\subsection{Uses}

Input Parameters Module

\subsection{Syntax}

\subsubsection{Exported Constants}

\begin{itemize}
    \item $\mathcal{D}_{\mathrm{MEG}}$: the collection of MEG recordings and associated metadata available to \progname.
\end{itemize}

\subsubsection{Exported Access Programs}

\begin{center}
\begin{tabular}{p{3.2cm} p{4.2cm} p{4.2cm} p{2.4cm}}
\hline
\textbf{Name} & \textbf{In} & \textbf{Out} & \textbf{Exceptions} \\
\hline
initMEGData & - & - & unsupported dataset, file error \\
getSubjectList & - & sequence of strings & dataset not initialized \\
getMEGResponse & subject id: string, utterance id: string & $\mathbf{Y}$ & subject not found, trial not found, file error \\
getSensorCount & - & $C$ & dataset not initialized \\
getMEGSamplingRate & - & $f_s$ & dataset not initialized \\
getTrialMetadata & subject id: string, utterance id: string & tuple & subject not found, trial not found \\
getStimulusLink & subject id: string, utterance id: string & string & subject not found, trial not found \\
\hline
\end{tabular}
\end{center}

\subsection{Semantics}

\subsubsection{State Variables}

\begin{itemize}
    \item $subjectIndex$: a mapping from subject identifiers to MEG recording descriptors.
    \item $trialIndex$: a mapping from $(subject\ id,\ utterance\ id)$ pairs to dataset locations (path) and metadata (wav file for audio, and bid-format MEG data for brain signal).
    \item $sensorInfo$: metadata describing sensor identities and count $C$.
    \item $initialized$: $\mathbb{B}$, indicates whether the module has loaded the MEG dataset structure.
\end{itemize}

\subsubsection{Environment Variables}

\begin{itemize}
    \item MEG recordings and metadata stored in the external file system.
    \item Dataset directory paths provided by the Input Parameters Module.
\end{itemize}

\subsubsection{Assumptions}

\begin{itemize}
    \item The MEG dataset is available at the configured location.
    \item Each pair of subject identifier and utterance identifier uniquely identifies one MEG trial.
    \item The stimulus linkage between MEG trials and utterance identifiers is valid and consistent.
\end{itemize}

\subsubsection{Access Routine Semantics}

\noindent \textbf{initMEGData}():
\begin{itemize}
    \item transition: $subjectIndex$, $trialIndex$, $sensorInfo$, and $initialized$ are constructed from the supplied MEG dataset configuration, and $initialized := \mathrm{true}$.
    \item exception: raised if the MEG dataset is unavailable, unsupported, or cannot be indexed.
\end{itemize}

\noindent \textbf{getSubjectList}():
\begin{itemize}
    \item output: returns the sequence of subject identifiers available in $\mathcal{D}_{\mathrm{MEG}}$.
    \item exception: raised if the module has not been initialized.
\end{itemize}

\noindent \textbf{getMEGResponse}(subject id, utterance id):
\begin{itemize}
    \item output: Read the preprocessed MEG trial from the external file system, and returns the response matrix $\mathbf{Y} \in \mathbb{R}^{N \times C}$, where $N$ is the number of time samples in the trial and $C$ is the number of MEG sensors defined by $sensorInfo$.
    \item exception: raised if the subject identifier or utterance identifier is unknown, or if the corresponding recording cannot be read.
\end{itemize}

\noindent \textbf{getSensorCount}():
\begin{itemize}
    \item output: looks up dataset in $sensorInfo$ and returns the sensor count $C$ stored in the corresponding sensor descriptor.
    \item exception: raised if the module has not been initialized.
\end{itemize}

\noindent \textbf{getMEGSamplingRate}():
\begin{itemize}
    \item output: returns the MEG sampling frequency $f_s$.
    \item exception: raised if the module has not been initialized.
\end{itemize}

\noindent \textbf{getTrialMetadata}(subject id, utterance id):
\begin{itemize}
    \item output: returns metadata associated with the given trial, such as timing information, recording identifiers, and preprocessing descriptors.
    \item exception: raised if the subject identifier or utterance identifier is unknown.
\end{itemize}

\noindent \textbf{getStimulusLink}(subject id, utterance id):
\begin{itemize}
    \item output: returns the stimulus identifier linked to the requested MEG trial.
    \item exception: raised if the subject identifier or utterance identifier is unknown.
\end{itemize}












\newpage




\section{MIS of Input Parameters Module} \label{Module_InputParameters}

\subsection{Module}

Input Parameters Module

\subsection{Uses}

None

\subsection{Syntax}

\subsubsection{Exported Constants}

\begin{itemize}
    \item $SupportedDatasets$: the set of supported dataset identifiers.
    \item $SupportedTasks$: the set of supported task identifiers.
    \item $SupportedLabelTypes$: the set of supported label-type identifiers.
    \item $SupportedPolicies$: the set of supported policy identifiers.
\end{itemize}

\subsubsection{Exported Access Programs}

\begin{center}
\begin{tabular}{p{3.4cm} p{4.8cm} p{4.0cm} p{3.0cm}}
\hline
\textbf{Name} & \textbf{In} & \textbf{Out} & \textbf{Exceptions} \\
\hline
getPreset &
preset name: string, seed: int &
TrainConfig &
unknown preset \\

validateConfig &
cfg: TrainConfig &
-- &
invalid numeric range, invalid categorical value, invalid field combination, file not found \\

ensureDirs &
cfg: TrainConfig &
-- &
directory creation failure \\

runId &
cfg: TrainConfig &
string &
-- \\

getRunDir &
cfg: TrainConfig &
Path &
-- \\

getCheckpointDir &
cfg: TrainConfig &
Path &
-- \\

getLogDir &
cfg: TrainConfig &
Path &
-- \\
\hline
\end{tabular}
\end{center}

\subsection{Semantics}

\subsubsection{State Variables}

\begin{itemize}
    \item $presetTable$: a mapping from preset names to preset-construction routines.
    \item $cfg$: the current training configuration object.
\end{itemize}

\subsubsection{Environment Variables}

\begin{itemize}
    \item Optional local path constants imported from the local path-definition file.
\end{itemize}

\subsubsection{Assumptions}

\begin{itemize}
    \item A requested preset name is intended to identify one preset-construction routine in $presetTable$.
    \item When path checking is enabled, the required dataset paths and manifest paths exist in the external file system.
\end{itemize}

\subsubsection{Access Routine Semantics}

\noindent \textbf{getPreset}(preset name, seed):
\begin{itemize}
    \item output: looks up the requested preset name in $presetTable$, invokes the corresponding preset-construction routine with the supplied seed, and returns a \texttt{TrainConfig}.
    \item exception: raised if the preset name is unknown.
\end{itemize}

\noindent \textbf{validateConfig}(cfg):
\begin{itemize}
    \item output: checks that the configuration satisfies the required numeric, categorical, and combination constraints, and optionally checks dataset paths.
    \item exception: raised if any configuration field violates a range
    constraint, categorical constraint, combination constraint, or path-validity requirement.
\end{itemize}

\noindent \textbf{ensureDirs}(cfg):
\begin{itemize}
    \item output: creates the run directory, checkpoint directory, and log
    directory associated with the supplied configuration.
    \item exception: raised if any required directory cannot be created.
\end{itemize}

\noindent \textbf{runId}(cfg):
\begin{itemize}
    \item output: returns the canonical run identifier derived from dataset, task, label type, policy, and seed.
\end{itemize}

\noindent \textbf{getRunDir}(cfg):
\begin{itemize}
    \item output: returns the run directory path associated with the supplied configuration.
\end{itemize}

\noindent \textbf{getCheckpointDir}(cfg):
\begin{itemize}
    \item output: returns the checkpoint directory path associated with the supplied configuration.
\end{itemize}

\noindent \textbf{getLogDir}(cfg):
\begin{itemize}
    \item output: returns the log directory path associated with the supplied configuration.
\end{itemize}

\subsubsection{Local Functions}

\begin{itemize}
    \item \textbf{\_coerce\_enum}: converts user-supplied values into the required enum type.
    \item \textbf{\_validate\_paths}: checks dataset and manifest paths when path checking is enabled.
    \item \textbf{\_apply\_task\_defaults}: fills task-specific default values in the configuration object.
\end{itemize}
















\newpage




\section{MIS of Model Training Module} \label{Module_ModelTraining}

\subsection{Module}

Model Training Module

\subsection{Uses}

Input Parameters Module, Audio Data Module

\subsection{Syntax}

\subsubsection{Exported Constants}

\begin{itemize}
    \item $\mathcal{C}_{\mathrm{Train}}$: the abstract type of a training configuration.
    \item $\mathcal{R}_{\mathrm{Eval}}$: the abstract type of an evaluation result record.
\end{itemize}

\subsubsection{Exported Access Programs}

\begin{center}
\begin{tabular}{p{3.4cm} p{4.8cm} p{4.0cm} p{3.0cm}}
\hline
\textbf{Name} & \textbf{In} & \textbf{Out} & \textbf{Exceptions} \\
\hline
getTrainer &
cfg: $\mathcal{C}_{\mathrm{Train}}$ &
training routine &
unsupported task or policy \\

runOnce &
cfg: $\mathcal{C}_{\mathrm{Train}}$, device: string &
result summary file and checkpoint files &
invalid configuration, loader failure, training failure, file error \\

buildLoaders &
cfg: $\mathcal{C}_{\mathrm{Train}}$, vocab, tailFrames: int &
train loader, validation loader, test loader &
dataset unavailable, malformed manifest, loader construction failure \\

runEval &
model, loader, device: string, srvTable, cfg: $\mathcal{C}_{\mathrm{Train}}$, tailFrames: int &
$\mathcal{R}_{\mathrm{Eval}}$ &
evaluation failure, malformed batch \\

saveCheckpoint &
path: Path, model, optimizer, epoch: int, bestVal: real &
-- &
file error \\
\hline
\end{tabular}
\end{center}

\subsection{Semantics}

\subsubsection{State Variables}

\begin{itemize}
    \item $cfg$: the active training configuration.
    \item $bestVal$: the best validation loss observed so far.
    \item $currentEpoch$: the current training epoch.
\end{itemize}

\subsubsection{Environment Variables}

\begin{itemize}
    \item Checkpoint files, logs, and result files stored in the external file system.
    \item The hardware device selected for training or evaluation (for example, CPU or CUDA device).
\end{itemize}

\subsubsection{Assumptions}

\begin{itemize}
    \item The supplied configuration has already been validated by the Input Parameters Module.
    \item The requested training dataset and manifests are available through the Audio Data Module.
    \item The selected model architecture is compatible with the batch feature dimension produced by the Audio Data Module.
\end{itemize}

\subsubsection{Access Routine Semantics}

\noindent \textbf{getTrainer}(cfg):
\begin{itemize}
    \item output: returns the training routine associated with the task and policy specified in $cfg$.
    \item exception: raised if no supported training routine exists for the requested configuration.
\end{itemize}

\noindent \textbf{runOnce}(cfg, device):
\begin{itemize}
    \item transition: creates output directories, initializes logging, constructs the vocabulary and data loaders, initializes the model and optimizer, runs the training loop, evaluates the best checkpoint on the test split, and writes checkpoint and result files.
    \item output: produces checkpoint files, logs, and a result summary for the completed training run.
    \item exception: raised if configuration validation, loader construction, training, evaluation, or output serialization fails.
\end{itemize}

\noindent \textbf{buildLoaders}(cfg, vocab, tailFrames):
\begin{itemize}
    \item output: returns the train, validation, and test data loaders implied by the current configuration and vocabulary.
    \item exception: raised if the required manifests or dataset subsets are unavailable or malformed.
\end{itemize}

\noindent \textbf{runEval}(model, loader, device, srvTable, cfg, tailFrames):
\begin{itemize}
    \item output: evaluates the model on the supplied loader and returns an evaluation-result record containing loss and summary metrics.
    \item exception: raised if evaluation encounters malformed inputs or invalid tensor shapes.
\end{itemize}

\noindent \textbf{saveCheckpoint}(path, model, optimizer, epoch, bestVal):
\begin{itemize}
    \item output: writes a checkpoint file containing model state, optimizer state, epoch number, and the best validation value observed so far.
    \item exception: raised if the checkpoint file cannot be written.
\end{itemize}

\subsubsection{Local Functions}

\begin{itemize}
    \item \textbf{srvLossAndCosFrame}: computes the frame-level SRV loss and cosine summary used during training and evaluation.
    \item \textbf{srvLossAndCosWord}: computes the word-level semantic SRV loss and cosine summary used during training and evaluation.
\end{itemize}












\newpage





\section{MIS of Frozen Model Module} \label{Module_FrozenModel}

\subsection{Module}

Frozen Model Module

\subsection{Uses}

Input Parameters Module, Model Training Module

\subsection{Syntax}

\subsubsection{Exported Constants}

\begin{itemize}
    \item $\mathcal{M}_{\mathrm{Frozen}}$: the abstract type of a frozen neural model.
\end{itemize}

\subsubsection{Exported Access Programs}

\begin{center}
\begin{tabular}{p{3.6cm} p{4.8cm} p{3.8cm} p{3.0cm}}
\hline
\textbf{Name} & \textbf{In} & \textbf{Out} & \textbf{Exceptions} \\
\hline
loadCheckpoint\\Model &
checkpoint path: Path, device: string &
$\mathcal{M}_{\mathrm{Frozen}}$ &
checkpoint unavailable, incompatible checkpoint, file error \\

freezeModel &
model &
$\mathcal{M}_{\mathrm{Frozen}}$ &
invalid model instance \\

getLayerCount &
frozen model: $\mathcal{M}_{\mathrm{Frozen}}$ &
$L$ &
invalid model instance \\

isFrozen &
frozen model: $\mathcal{M}_{\mathrm{Frozen}}$ &
$\mathbb{B}$ &
invalid model instance \\
\hline
\end{tabular}
\end{center}

\subsection{Semantics}

\subsubsection{State Variables}

\begin{itemize}
    \item $loadedCheckpoint$: the checkpoint currently loaded by the module.
    \item $frozenModel$: the non-trainable model instance derived from the loaded checkpoint.
\end{itemize}

\subsubsection{Environment Variables}

\begin{itemize}
    \item Stored checkpoint files in the external file system.
\end{itemize}

\subsubsection{Assumptions}

\begin{itemize}
    \item A valid trained checkpoint exists at the supplied checkpoint path.
    \item The checkpoint contains sufficient information to reconstruct the model architecture and load its parameters.
\end{itemize}

\subsubsection{Access Routine Semantics}

\noindent \textbf{loadCheckpointModel}(checkpoint path, device):
\begin{itemize}
    \item transition: reads the checkpoint file, reconstructs the corresponding model instance, loads the stored parameters, converts the model to evaluation mode, and stores it as $frozenModel$.
    \item output: returns the frozen model instance.
    \item exception: raised if the checkpoint is unavailable, incompatible, or cannot be loaded.
\end{itemize}

\noindent \textbf{freezeModel}(model):
\begin{itemize}
    \item output: disables parameter updates for the supplied model and returns the frozen model instance.
    \item exception: raised if the supplied model is invalid.
\end{itemize}

\noindent \textbf{getLayerCount}(frozen model):
\begin{itemize}
    \item output: returns the number of accessible internal layers exposed by the frozen model.
    \item exception: raised if the supplied model is invalid.
\end{itemize}

\noindent \textbf{isFrozen}(frozen model):
\begin{itemize}
    \item output: returns \texttt{true} if the supplied model is in non-trainable frozen form, and \texttt{false} otherwise.
    \item exception: raised if the supplied model is invalid.
\end{itemize}

\subsubsection{Local Functions}

\begin{itemize}
    \item \textbf{inferModelFromCheckpoint}: reconstructs the model architecture associated with a stored checkpoint.
    \item \textbf{loadStateDict}: loads stored parameters into the reconstructed model.
\end{itemize}













\newpage






\section{MIS of Representation Extraction Module} \label{Module_RepresentationExtraction}

\subsection{Module}

Representation Extraction Module

\subsection{Uses}

Input Parameters Module, Audio Data Module, Frozen Model Module, MEG Data Module

\subsection{Syntax}

\subsubsection{Exported Constants}

\begin{itemize}
    \item $\mathcal{H}$: the abstract type of a hidden-state sequence.
    \item $\mathcal{R}$: the abstract type of a time-by-feature predictor matrix.
\end{itemize}

\subsubsection{Exported Access Programs}

\begin{center}
\begin{tabular}{p{3.8cm} p{4.8cm} p{3.8cm} p{3.0cm}}
\hline
\textbf{Name} & \textbf{In} & \textbf{Out} & \textbf{Exceptions} \\
\hline
forwardWithHidden &
frozen model, input tensor, input lengths &
hidden-state sequence $\in \mathcal{H}$ &
model execution failed, invalid input shape \\

computePredictors\\FromHidden &
hidden-state sequence $\in \mathcal{H}$, transform specification &
predictor array(s) &
invalid hidden-state shape, transform failed \\

alignToAnalysisGrid &
predictor array(s), analysis time grid &
predictor matrix $\in \mathcal{R}$ &
time-grid mismatch, alignment failed \\

serializePredictors &
predictor matrix $\in \mathcal{R}$, output path specification &
saved predictor file(s) &
file error, serialization failed \\
\hline
\end{tabular}
\end{center}

\subsection{Semantics}

\subsubsection{State Variables}

\begin{itemize}
    \item $extractionSpec$: the extraction and transformation parameters supplied by the Input Parameters Module.
    \item $timeGrid$: the analysis time grid used to align predictor outputs.
\end{itemize}

\subsubsection{Environment Variables}

\begin{itemize}
    \item Input feature files and serialized predictor files stored in the external file system.
\end{itemize}

\subsubsection{Assumptions}

\begin{itemize}
    \item The required frozen model has already been loaded.
    \item The requested input sequence has a feature dimension compatible with the frozen model.
    \item The analysis time grid is available and compatible with the predictor length.
\end{itemize}

\subsubsection{Access Routine Semantics}

\noindent \textbf{forwardWithHidden}(frozen model, input tensor, input lengths):
\begin{itemize}
    \item output: applies the frozen model to the supplied input sequence and returns the hidden-state sequence exposed by the model.
    \item exception: raised if the model is unavailable or the input tensor has an invalid shape.
\end{itemize}

\noindent \textbf{computePredictorsFromHidden}(hidden-state sequence, transform specification):
\begin{itemize}
    \item output: applies the requested transformation to the hidden-state sequence and returns one or more predictor arrays derived from the hidden representations.
    \item exception: raised if the hidden-state sequence is invalid or the requested transformation cannot be applied.
\end{itemize}

\noindent \textbf{alignToAnalysisGrid}(predictor array(s), analysis time grid):
\begin{itemize}
    \item output: converts the predictor array(s) into predictor matrices aligned to the requested analysis time grid.
    \item exception: raised if the predictor length and the analysis time grid are incompatible.
\end{itemize}

\noindent \textbf{serializePredictors}(predictor matrix, output path specification):
\begin{itemize}
    \item output: writes the aligned predictor matrix to the required serialized output format.
    \item exception: raised if serialization fails or the output files cannot be written.
\end{itemize}

\subsubsection{Local Functions}

\begin{itemize}
    \item \textbf{computeMagnitudeChange}: computes magnitude-based and change-based predictor arrays from hidden states.
    \item \textbf{wrapToSerializedPredictor}: converts predictor arrays into the serialized predictor format expected by downstream analysis.
\end{itemize}











\newpage



\section{MIS of Statistical Analysis Module} \label{Module_StatisticalAnalysis}

\subsection{Module}

Statistical Analysis Module

\subsection{Uses}

Input Parameters Module, MEG Data Module, Representation Extraction Module

\subsection{Syntax}

\subsubsection{Exported Constants}

\begin{itemize}
    \item $\mathcal{T}_{\mathrm{TRF}}$: the abstract type of a TRF analysis result.
    \item $\mathcal{S}_{\mathrm{Summary}}$: the abstract type of a statistical summary table.
\end{itemize}

\subsubsection{Exported Access Programs}

\begin{center}
\begin{tabular}{p{3.8cm} p{4.8cm} p{3.8cm} p{3.0cm}}
\hline
\textbf{Name} & \textbf{In} & \textbf{Out} & \textbf{Exceptions} \\
\hline
runTRFJob &
analysis configuration &
$\mathcal{T}_{\mathrm{TRF}}$ &
invalid configuration, predictor unavailable, fitting failed \\

loadPredictorSafe &
predictor identifier, predictor source &
predictor object &
predictor unavailable, non-finite values, malformed predictor \\

summarizeTRF\\Results &
TRF result $\in \mathcal{T}_{\mathrm{TRF}}$ &
$\mathcal{S}_{\mathrm{Summary}}$ &
result unavailable, malformed result \\

computeImprovement\\Table &
baseline result, target result &
improvement summary table &
shape mismatch, comparison failed \\

selectBestLayer &
summary table, selection criterion &
layer index &
invalid criterion, empty summary \\
\hline
\end{tabular}
\end{center}

\subsection{Semantics}

\subsubsection{State Variables}

\begin{itemize}
    \item $analysisSpec$: the statistical-analysis configuration supplied by the Input Parameters Module.
    \item $resultCache$: a mapping from analysis identifiers to stored TRF result objects.
\end{itemize}

\subsubsection{Environment Variables}

\begin{itemize}
    \item Cached TRF result files stored in the external file system.
\end{itemize}

\subsubsection{Assumptions}

\begin{itemize}
    \item The required MEG data and aligned predictors are already available.
    \item The analysis configuration defines a valid lag window, sampling rate, and cross-validation setting.
    \item The underlying Eelbrain/TRF pipeline is available to execute the requested analysis.
\end{itemize}

\subsubsection{Access Routine Semantics}

\noindent \textbf{runTRFJob}(analysis configuration):
\begin{itemize}
    \item transition: executes the configured TRF analysis, stores the resulting TRF object in $resultCache$, and records any generated cache files.
    \item output: returns the TRF result object for the requested analysis job.
    \item exception: raised if the supplied analysis configuration is invalid or if TRF execution fails.
\end{itemize}

\noindent \textbf{loadPredictorSafe}(predictor identifier, predictor source):
\begin{itemize}
    \item output: loads the requested predictor, checks for non-finite values and degenerate columns, and returns a predictor object suitable for downstream TRF analysis.
    \item exception: raised if the predictor is unavailable, malformed, or unusable.
\end{itemize}

\noindent \textbf{summarizeTRFResults}(TRF result):
\begin{itemize}
    \item output: extracts the summary quantities required for downstream reporting, comparison, and plotting.
    \item exception: raised if the supplied TRF result is malformed or incomplete.
\end{itemize}

\noindent \textbf{computeImprovementTable}(baseline result, target result):
\begin{itemize}
    \item output: returns the baseline-referenced comparison table between the supplied analysis results.
    \item exception: raised if the supplied results are incompatible.
\end{itemize}

\noindent \textbf{selectBestLayer}(summary table, selection criterion):
\begin{itemize}
    \item output: returns the layer index that optimizes the supplied selection criterion.
    \item exception: raised if the summary table is empty or the criterion is invalid.
\end{itemize}

\subsubsection{Local Functions}

\begin{itemize}
    \item \textbf{labelEvents}: matches event markers to stimulus structure and constructs the event labels used for epoching.
    \item \textbf{sanitizePredictor}: replaces non-finite predictor values and detects flat columns before analysis.
\end{itemize}


















\newpage



\section{MIS of Plotting Module} \label{Module_Plotting}

\subsection{Module}

Plotting Module

\subsection{Uses}

Input Parameters Module, Statistical Analysis Module

\subsection{Syntax}

\subsubsection{Exported Constants}

None

\subsubsection{Exported Access Programs}

\begin{center}
\begin{tabular}{p{3.6cm} p{4.8cm} p{3.8cm} p{3.0cm}}
\hline
\textbf{Name} & \textbf{In} & \textbf{Out} & \textbf{Exceptions} \\
\hline
plotCurves &
cosine curves, probability curves, $L$, env sampling rate, output path, title &
saved figure path &
invalid plotting input, plotting failed, file error \\

buildPlotSpec &
analysis summary, plotting options &
plot specification &
invalid summary, invalid plotting options \\

saveFigure &
plot specification, output path &
saved figure path &
plotting failed, file error \\
\hline
\end{tabular}
\end{center}

\subsection{Semantics}

\subsubsection{State Variables}

\begin{itemize}
    \item $figureRegistry$: a mapping from plot identifiers to saved figure paths.
\end{itemize}

\subsubsection{Environment Variables}

\begin{itemize}
    \item Figure files written to the external file system.
\end{itemize}

\subsubsection{Assumptions}

\begin{itemize}
    \item The statistical summaries or curve arrays required for plotting are already available.
    \item The output directory for the requested figure is writable.
\end{itemize}

\subsubsection{Access Routine Semantics}

\noindent \textbf{plotCurves}(cosine curves, probability curves, $L$, env sampling rate, output path, title):
\begin{itemize}
    \item transition: renders the supplied curve arrays into a figure and stores the saved figure path in $figureRegistry$.
    \item output: returns the path of the generated figure file.
    \item exception: raised if the supplied curve arrays are malformed or the figure cannot be written.
\end{itemize}

\noindent \textbf{buildPlotSpec}(analysis summary, plotting options):
\begin{itemize}
    \item output: returns the internal plotting specification derived from the supplied analysis summary and plotting options.
    \item exception: raised if the summary or plotting options are invalid.
\end{itemize}

\noindent \textbf{saveFigure}(plot specification, output path):
\begin{itemize}
    \item transition: renders the supplied plotting specification and writes the corresponding figure file to the requested output path.
    \item output: returns the saved figure path.
    \item exception: raised if rendering fails or the output path is invalid.
\end{itemize}

\subsubsection{Local Functions}

\begin{itemize}
    \item \textbf{renderFigure}: converts the plotting specification into a concrete figure object.
    \item \textbf{registerFigure}: records the saved figure path in $figureRegistry$.
\end{itemize}







\bibliographystyle {plainnat}
\bibliography {../../../refs/References}















\newpage

% \section{Appendix} \label{Appendix}

% \wss{Extra information if required}

% \newpage{}

% \section*{Appendix --- Reflection}

% \wss{Not required for CAS 741 projects}

% The information in this section will be used to evaluate the team members on the
% graduate attribute of Problem Analysis and Design.

% \input{../../Reflection.text}

% \begin{enumerate}
%   \item What went well while writing this deliverable? 
%   \item What pain points did you experience during this deliverable, and how
%     did you resolve them?
%   \item Which of your design decisions stemmed from speaking to your client(s)
%   or a proxy (e.g. your peers, stakeholders, potential users)? For those that
%   were not, why, and where did they come from?
%   \item While creating the design doc, what parts of your other documents (e.g.
%   requirements, hazard analysis, etc), it any, needed to be changed, and why?
%   \item What are the limitations of your solution?  Put another way, given
%   unlimited resources, what could you do to make the project better? (LO\_ProbSolutions)
%   \item Give a brief overview of other design solutions you considered.  What
%   are the benefits and tradeoffs of those other designs compared with the chosen
%   design?  From all the potential options, why did you select the documented design?
%   (LO\_Explores)
%   \item (After you have implemented another team's module, which means this
%   isn't filled in until after the original deadline). What did you learn by
%   implementing another team's module? Were all the details you needed in the
%   documentation, or did you need to make assumptions, or ask the other team
%   questions?  If your team also had another team implement one of your modules,
%   what was this experience like?  Are there things in your documentation you
%   could have changed to make the process go more smoothly for when an
%   ``outsider'' completes some of the implementation?
% \end{enumerate}


\end{document}